\Lecturer{Hui Jin}
\Contributor{}
\LectureNumber{07}
\LectureDate{DATE}
\LectureTitle{Convolution Codes - Viterbi Decoding}

\lstset{style=mystyle}

	\MakeScribeTop

%#############################################################
%#############################################################
%#############################################################
%#############################################################

\section{Introduction}
If bit streams are encoded by a convolutional code, how do we find the correct codeword \(c\) given a received vector \(y\) which is a noisy version of a codeword? As noises are stochastic, finding the correct codeword is not guaranteed; instead our task is to find the most likely codeword.

For convolutional coding, there exists a low complexity  decoding algorithm, called Viterbi decoding, that guarantees to find the maximum likelihood codeword. It should be noted that even though Dr. Andrew Viterbi proposed the algorithm in 1967, it was not until six years later in 1973 that David Forney proved its optimality. 

In this lecture we shall describe this amazing algorithm. We will learn that Viterbi decoding is really a dynamic programming algorithm applied to a Markov tree process; the Markov tree corresponds to the trellis structure of the convolutional code.   

But before we discuss the Viterbi decoding, we describe a hard decision decoder. On the BSC channel, it is the optimal decoder. On the AWGN channel, it is a simplied version of the optimal Viterbi decoder.




\section{Maximum Likelihood Decoder}
The quantity \(p(y|c)\) represents the probability of receiving the vector \(y\) given that codeword \(c\) is transmitted. A Maximum Likelihood Decoder chooses the codeword such that this probability is maximized. 

If we assume the channel is memoryless, then for each transmitted bit, the noise generated by the channel is independent. That is to say
\begin{equation}
p(y|c) = \prod_{i=1}^{n}p(y_i|c_i) \label{eq:pyc}
\end{equation}
where \(c_i\) represents the block of parity bits transmitted at the \(i\)th input block, and \(y_i\) represents the block of received value after the channel if \(c_i\) is transmitted. Note that each block of parity bits \(c_i\) itself is a vector, consisting of \(r\) bits, i.e. \(c_i =c_{i}^{1}c_{i}^{1}\cdots c_{i}^{r}.\) Similarly, the block of the received value is a vector of length 
\(r\) as well: \(y_i = y_{i}^{1}y_{i}^{2}\cdots y_{i}^{r}\). Even though Eq.~\ref{eq:pyc} can be further factored into as \(p(y_i|c_i) = \prod_{j=1}^{r}p(y_i^{j}|c_i^{j})\), we hold off using this lest of unnecessary cluttering. We will use it in Section 3. 


Our maximum likelihood decoder chooses the codeword \(c\) such that the probability \(p(y|c)\) is maximized. 
\begin{equation}
    c_{ML} = \argmax_{c} p(y|c) = \argmax_{c=c_1, c_2, \cdots, c_n}  \prod_{i=1}^{n}p(y_i|c_i) \label{eq:MLDecoder}
\end{equation}

An intuitive understanding of the ML decoder is: on the trellis we want to find the path, which corresponds to a codeword, such that the probability of that path given received vector \(y\) is the largest. The decoding capability of convolutional codes comes from the fact that the distances between valid paths can be large. This is essentially the free distance property of the convolutional codes.  

\section{A hard decision decoder}


\section{A soft decision Decoder}
\subsection{Log-likelihood}
Since logrithm is a monotonical function, Eq~(\ref{eq:MLDecoder}) is equivalent to 
\begin{align}
    c_{ML} &= \argmax_{c} \log p(y|c) \nonumber \\ 
           &= \argmax_{c=c_1, c_2, \cdots, c_n} \log \prod_{i=1}^{n}p(y_i|c_i) \nonumber \\
           &= \argmax_{c=c_1, c_2, \cdots, c_n} \sum_{i=1}^{n} \log p(y_i|c_i) \label{eq:ML_Viterbi}
\end{align}


\subsection{Normalization and Log-likelihood Ratio}
We also realize that subtracting a constant from each term in Eq.~\ref{eq:ML_Viterbi} does not change the argmax. A neat trick to simplify the decoder is to subtract \(\log p(y_i|\bf{1})\) in the \(i\)-th sum, here \(\log p(y_i|\bf{1})\) represents the log-likeliood of receiving \(y_i\) given the all-one vector is transmitted in the \(i\)-th block. (It does not matter we use which transmitted vector as the base; we choose all one vector due to its simplicity in description.) 

\begin{align}
    c_{ML} &= \argmax_{c=c_1, c_2, \cdots, c_n} \sum_{i=1}^{n} (\log p(y_i|c_i) - \log p(y_i|0) \nonumber ) \\
    &=\argmax_{c=c_1, c_2, \cdots, c_n} \sum_{i=1}^{n} \log \frac{p(y_i|c_i)}{p(y_i|\bf{1})} \label{eq:CovnCML}
\end{align}

The quantity \( \log \frac{p(y_i|c_i)}{p(y_i|\bf{1})} \) is known the log-likelihood ratio between any other branch and the branch that has the all one transmitted bits. This trick is useful to remove the constant in  \(p(y_i|c_i)\); and it also nicely limits the data range in actual software/hardware implementation.

\begin{example}
On the Binary Symmetric Channel, assuming the two transmitted parity bits are \(c_i = 01\), the received vector is \(y_i = 11\), then we have 
\[p(y_i|c_i) = p(1-p),\]
and the log-likelihood is 
\[\log p(y_i|c_i) = \log(p(1-p)),\]
and the normalized log-likelihood ratio is

\begin{align*}
&\log \frac{p(y_i|c_i)}{p(y_i|11)}) \\ &= \log \frac{p(1-p)}{(1-p)(1-p)} \\
&=\log \frac{p}{1-p}.
\end{align*}

\end{example}

\begin{example}
On the AWGN Channel with Gaussian noise \(N(0, \sigma^2)\), assuming the two transmitted parity bits are \(c_i = 01\), the received vector is \(y_i = y_0y_1\), then we have
\[p(y_i|c_i) = \frac{1}{2\pi\sigma^2} \exp[-\frac{(y_0-1)^2+(y_1 + 1)^2}{2\sigma^2}],\]
and the log-likelihood is 
\[\log p(y_i|c_i) = \log(\frac{1}{2\pi\sigma^2}) -\frac{(y_0-1)^2+(y_1 + 1)^2}{2\sigma^2} ,\]
and the normalized log-likehood ratio is
\begin{align*}
&\log \frac{p(y_i|c_i)}{p(y_i|11)})\\ &= -\frac{(y_0-1)^2+(y_1 + 1)^2}{2\sigma^2} + \frac{(y_0+1)^2+(y_1 + 1)^2}{2\sigma^2} \\ 
&= \frac{2y_0}{\sigma^2}
\end{align*}
This last expression is surprisingly simple. And it makes intuitive sense. A positive \(y_0\) indicates the transmitted bit \(c_0\) is more likely to be \(0\), as the branch metric is \(\log \frac{p(y_0|0)}{p(y_0|1)}\). A negative says otherwise. 
\end{example}



\section{Branch Metric}
Let us define the \(i\)th branch metric to be 
\begin{equation}
  B_i(c_i) = \log \frac{p(y_i|c_i)}{p(y_i|\bf{1})} \label{eq:BranchMetric}
\end{equation}

Therefore, the maximum likelihood decision in Eq.~\ref{eq:CovnCML} turns to:
\begin{align}
    c_{ML}
    &=\argmax_{c=c_1, c_2, \cdots, c_n} \sum_{i=1}^{n} \log B_i(c_i) \label{eq:ConvBranchML} 
\end{align}
We will stamp the branch metric on the trellis and now the trellis graph has weights between the \(i\)-th state and the \((i+1)\)-state, for any pair of valid states. 

Let us look at two examples, both the binary symmetric channel case and the AWGN channel case. 

The maximum likelihood decoding problems becomes the problem of finding the maximum weight path along this trellis tree.   

Shed in this light, the optimal decision problem changes to a well known approach - dynamic programming! 

\section{Dynamic Programming}
Let us denote the trellis states as \(0, 1, \ldots, 2^K-1\), where \(K\) is the constraint length. In our example, \(K=2\), and the states are \(0, 1, 2, 3\).

The maximum likelihood decoded codeword corresponds to the path from state \(0\) at time \(t=1\) to any final state \(s\) at time \(t=n\) with the maximum sum of branch metrics. If we find the maximum path, then the associated path correspond to the Maximum Likelihood codeword. So our problem is really about the maximum path branch: 
\begin{align}
    B_{ML} &=\max_{c=c_1, c_2, \cdots, c_n} \sum_{i=1}^{n} \log B_i(c_i) \label{eq:ConvBranchML} 
\end{align}

But the maximum sum of branch metrics at time \(t=n\) can be decomposed to the maximum sum of branch metrics at time \(t=n-1\) and the branch metric at time \(t=n\). That is the basic principle of dynamic programming.

Formally, let us define the sum of a partial path, from state \(0\) at time \(t=1\) to a state \(s\) at time \(t=j\) as the cumulative branch metric (CB):
\begin{equation*}
    CB_{j}(s) = \max_{c=c_1, c_2, \cdots, c_j }\sum_{i=1}^{n} \log B_i(c_i),
\end{equation*}
where the ending state of \(c_1, c_2, \cdots, c_j\) is \(s\).

Clearly \(B_{ML} = \max_s CB_{n}(s).\) To compute \(CB_{j}(s)\), we rely on the following insight: 
\begin{theorem}
For a state \(s\), let \(P(s)\) be the set of states that can transit to \(s\). For state \(c' \in Pre(s)\), let \(c(s',s)\) be the coded bits generated as state \(s'\) transit to \(s\). With that definition, we have  
\begin{equation*}
    CB_j(s) = \max_{s' \in P(s)} (CB_{j-1}(s') + B_j(c(s',s))
\end{equation*}    
\end{theorem}
\begin{proof}
    It is almost evident if we look at the trellis. 
\end{proof}

\section{Put it together - the Viterbi Algorithm}
As described, the Viterbi algorithm is straightforward: 
\begin{itemize}
    \item calculate the branch metrics
    \item forward path dynamic programming
    \item trace back the maximum path to get the codeword
\end{itemize}


We will look at the trellis structure in figure. 

The following examples will reinforce our understanding. Costello book? Wicker book?  

\begin{example}
    
\end{example}

\section{Varying the code rate}
We can achieve different code rate through puncturing. 


\section{The complexity of Viterbi Decoding}
It only requires a forward path. At each time step, we have \(2^K \times M\) computations, total operations is \(n M 2^K\). 

Storage requirement: about \(2^K\) branch metrics and record the codeword. 


\section{The performance of Convolution Codes}

The 3dB performance gap from the Shannon limit. 

